\section{Notation}
\label{sec:notation}

\subsection{Face turns}

A letter names a face and the letter alone means a quarter turn clockwise as seen
from outside that face.  The suffix \texttt{2} means a half turn and \texttt{'}
means a quarter turn anticlockwise (WCA Regulation 12a1).

\begin{center}
\begin{tabular}{@{}llll@{}}
\toprule
Letter & Face & Clockwise seen from & Opposite \\
\midrule
\texttt{U} & upper        & above   & \texttt{D} \\
\texttt{D} & bottom       & below   & \texttt{U} \\
\texttt{R} & right        & right   & \texttt{L} \\
\texttt{L} & left         & left    & \texttt{R} \\
\texttt{F} & front        & front   & \texttt{B} \\
\texttt{B} & back         & behind  & \texttt{F} \\
\bottomrule
\end{tabular}
\end{center}

There are 18 face turns: \texttt{U U2 U' R R2 R' F F2 F' D D2 D' L L2 L' B B2 B'}.

\subsection{Wide turns}

An \texttt{w} after the letter turns the outer face \emph{and} the layer behind it
(Regulation 12a2): \texttt{Uw Dw Lw Rw Fw Bw} with the usual suffixes, 18 moves in
total.  The layer count may be written out, so \texttt{2Rw} is the same move as
\texttt{Rw}; \texttt{1Rw} is not valid notation and \texttt{3Rw} does not exist on a
$3\times3\times3$.

\subsection{Rotations and slices}

A rotation turns the \emph{whole} cube, including its centres, in the direction of
a face turn (Regulation 12a4): \texttt{x} follows \texttt{R}, \texttt{y} follows
\texttt{U} and \texttt{z} follows \texttt{F}, with the usual suffixes.

The slice turns \texttt{M}, \texttt{E} and \texttt{S} turn only the middle layer,
following \texttt{L}, \texttt{D} and \texttt{F} respectively.  They are not part of
WCA notation but appear in many published algorithms, so the engine understands
them; Section~\ref{sec:tables} explains how the CLI keeps the display upright when
an algorithm uses them.

\subsection{Accepted aliases}

\begin{center}
\begin{tabular}{@{}ll@{}}
\toprule
Written & Read as \\
\midrule
\texttt{r l u d f b} & \texttt{Rw Lw Uw Dw Fw Bw} (older listings) \\
\texttt{m e s} & \texttt{M E S} \\
\texttt{RW} & \texttt{Rw} (the \texttt{w} is case insensitive) \\
\texttt{2Rw} & \texttt{Rw} \\
\texttt{R2'} or \texttt{R'2} & \texttt{R2} \\
\bottomrule
\end{tabular}
\end{center}

Anything else is rejected with the offending token and its position, for example
\texttt{error: layer prefix '3' is invalid on a 3x3x3 cube (index=4, token='3Rw')}.

\subsection{The facelet representation}
\label{sec:facelets}

The cube is a string of 54 colours in Kociemba's order: the nine facelets of
\texttt{U}, then of \texttt{R}, \texttt{F}, \texttt{D}, \texttt{L} and \texttt{B}.
Inside a face the facelets are numbered row by row in reading order \emph{as seen
from outside the cube}, with the following view directions, which are the ones
Kociemba's own move definitions use:

\begin{center}
\begin{tabular}{@{}lllll@{}}
\toprule
Face & View up & Top-left facelet belongs to & Top-right & Bottom-right \\
\midrule
\texttt{U} & back          & UBL & UBR & UFR \\
\texttt{R} & up            & UFR & UBR & DBR \\
\texttt{F} & up            & UFL & UFR & DFR \\
\texttt{D} & front         & DFL & DFR & DBR \\
\texttt{L} & up            & UBL & UFL & DFL \\
\texttt{B} & up            & UBR & UBL & DBL \\
\bottomrule
\end{tabular}
\end{center}

For example, \texttt{U9} is the upper sticker of the front-right corner,
\texttt{R1} is the right sticker of that same corner, and \texttt{F3} is its front
sticker, so the three facelets \texttt{(U9, R1, F3)} describe one and the same
corner piece.  The solved cube is

\begin{center}
\texttt{UUUUUUUUURRRRRRRRRFFFFFFFFFDDDDDDDDDLLLLLLLLLBBBBBBBBB}
\end{center}

and the six centre facelets --- \texttt{U5}, \texttt{R5}, \texttt{F5}, \texttt{D5},
\texttt{L5}, \texttt{B5} --- say which colour belongs to which face.  A state in
which the centres have moved is still a legal cube (it has simply been turned in
the hand); the program handles that by relabelling the colours before it reasons
about pieces.

\subsection{Moves as permutations}

Every move is stored as a permutation of the 54 facelets: for each facelet, where
its sticker goes.  The eighteen face turns were generated from the cube's geometry
and then checked, one by one, against the move tables published on
\texttt{kociemba.org}; wide turns are two-layer rotations, rotations are three-layer
rotations and slices are the single middle layer.  That is why a wide move or a
rotation can move a centre sticker and a face turn cannot.
